\section{Standard-halo integrals}\label{app:D.1}

This section evaluates the averages of the standard halo model used in \cref{sec:1.2,sec:1.3}: the normalization $N_{\mathrm{esc}}$, the laboratory speed distribution $F_\chi(v)$, the mean inverse speed $\langle 1/v\rangle$ and the function $\eta(\vmin)$ for every $\vmin$. The galactic-frame distribution is $f_{\mathrm{gal}}(\vect{v}_{\mathrm{gal}}) = N_{\mathrm{esc}}^{-1}\ee^{-v_{\mathrm{gal}}^2/v_0^2}\Theta(v_{\mathrm{esc}} - v_{\mathrm{gal}})$, and the laboratory distribution is $f_\chi(\vect{v}) = f_{\mathrm{gal}}(\vect{v} + \vect{v}_{\mathrm{E}})$ (\cref{eq:1.2-fgal,eq:1.2-flab}). Throughout, $v_{\mathrm{E}} = |\vect{v}_{\mathrm{E}}|$ at the time of interest, and we assume $0 < v_{\mathrm{E}} < v_{\mathrm{esc}}$; the values of \cref{app:A.8} give $v_{\mathrm{E}}/v_{\mathrm{esc}} \le 0.49$. Only one integral is needed:
\begin{equation}
  \int_{y_1}^{y_2} \ee^{-y^2/v_0^2}\,\dd y = \frac{\sqrt{\pi}\,v_0}{2}\left[\operatorname{erf}\!\left(\frac{y_2}{v_0}\right) - \operatorname{erf}\!\left(\frac{y_1}{v_0}\right)\right].
  \label{eq:D.1-erf}
\end{equation}
\Cref{eq:D.1-erf} is the definition of the error function (\cref{app:C.8}).

\subsection*{Normalization}

The integrand of $\int \dd^3 v_{\mathrm{gal}}\,\ee^{-v_{\mathrm{gal}}^2/v_0^2}\Theta(v_{\mathrm{esc}} - v_{\mathrm{gal}})$ depends only on $y = v_{\mathrm{gal}}$; in spherical coordinates the solid angle contributes $4\pi$ and $N_{\mathrm{esc}} = 4\pi\int_0^{v_{\mathrm{esc}}} y^2\ee^{-y^2/v_0^2}\dd y$. Integrating by parts with $y\,\ee^{-y^2/v_0^2}\dd y = \dd(-\tfrac12 v_0^2\ee^{-y^2/v_0^2})$,
\begin{multline*}
  \int_0^{v_{\mathrm{esc}}} y^2\ee^{-y^2/v_0^2}\dd y = \Bigl[-\tfrac12 v_0^2\,y\,\ee^{-y^2/v_0^2}\Bigr]_0^{v_{\mathrm{esc}}} + \tfrac12 v_0^2\int_0^{v_{\mathrm{esc}}}\ee^{-y^2/v_0^2}\dd y\\
  = -\tfrac12 v_0^2 v_{\mathrm{esc}}\ee^{-v_{\mathrm{esc}}^2/v_0^2} + \frac{\sqrt{\pi}\,v_0^3}{4}\operatorname{erf}\!\left(\frac{v_{\mathrm{esc}}}{v_0}\right).
\end{multline*}
Multiplying by $4\pi$ gives \cref{eq:1.2-nesc}.

\subsection*{Speed distribution}

The cut in \cref{eq:1.2-Fspeed} removes directions only for $v > v_{\mathrm{esc}} - v_{\mathrm{E}}$, because the galactic speed $|\vect{v} + \vect{v}_{\mathrm{E}}|$ ranges over $[\,|v - v_{\mathrm{E}}|, v + v_{\mathrm{E}}]$ as the direction of $\vect{v}$ varies. All directions survive if $v + v_{\mathrm{E}} < v_{\mathrm{esc}}$, none if $|v - v_{\mathrm{E}}| > v_{\mathrm{esc}}$, that is (since $v_{\mathrm{E}} < v_{\mathrm{esc}}$) if $v > v_{\mathrm{esc}} + v_{\mathrm{E}}$; in between, the surviving directions are $-1 \le \cos\vartheta < \cos\vartheta_{\mathrm{cut}}$ with $v^2 + v_{\mathrm{E}}^2 + 2vv_{\mathrm{E}}\cos\vartheta_{\mathrm{cut}} = v_{\mathrm{esc}}^2$. The angular integral is
\begin{equation*}
  \int_{-1}^{\cos\vartheta_{\mathrm{cut}}}\ee^{-2vv_{\mathrm{E}}\cos\vartheta/v_0^2}\,\dd(\cos\vartheta) = \frac{v_0^2}{2vv_{\mathrm{E}}}\left[\ee^{2vv_{\mathrm{E}}/v_0^2} - \ee^{-2vv_{\mathrm{E}}\cos\vartheta_{\mathrm{cut}}/v_0^2}\right].
\end{equation*}
Multiplying by $2\pi v^2 N_{\mathrm{esc}}^{-1}\ee^{-(v^2 + v_{\mathrm{E}}^2)/v_0^2}$ and completing the squares gives both branches of \cref{eq:1.2-Fspeed}, with prefactor $\pi v_0^2 v/(v_{\mathrm{E}}N_{\mathrm{esc}})$. We write it as $F_\chi(v) = K_{\mathrm{sp}}\,v\,[\ee^{-(v - v_{\mathrm{E}})^2/v_0^2} - e_{\mathrm{cut}}(v)]$, with $K_{\mathrm{sp}} = \pi v_0^2/(v_{\mathrm{E}}N_{\mathrm{esc}})$ and $e_{\mathrm{cut}}(v) = \ee^{-(v + v_{\mathrm{E}})^2/v_0^2}$ below $v_{\mathrm{esc}} - v_{\mathrm{E}}$, $e_{\mathrm{cut}}(v) = \ee^{-v_{\mathrm{esc}}^2/v_0^2}$ above it.

\subsection*{The function \texorpdfstring{$\eta(\vmin)$}{eta(vmin)}}

\Cref{sec:1.3} defines $\eta(\vmin) = \langle \Theta(v - \vmin)/v\rangle = \int_{\vmin}^\infty \dd v\,F_\chi(v)/v$. The factor $v$ in $F_\chi$ cancels the $1/v$, so for $\vmin \ge 0$
\begin{equation*}
  \eta(\vmin) = K_{\mathrm{sp}}\left[\int_{\vmin}^{v_{\mathrm{esc}} + v_{\mathrm{E}}}\ee^{-(v - v_{\mathrm{E}})^2/v_0^2}\dd v - \int_{\vmin}^{v_{\mathrm{esc}} + v_{\mathrm{E}}}e_{\mathrm{cut}}(v)\,\dd v\right],
\end{equation*}
where the upper limits are those of $F_\chi$.

\emph{(i) $0 \le \vmin < v_{\mathrm{esc}} - v_{\mathrm{E}}$.} With $y = v - v_{\mathrm{E}}$ the first integral is over $\vmin - v_{\mathrm{E}} < y < v_{\mathrm{esc}}$. The second has two pieces: $\ee^{-(v + v_{\mathrm{E}})^2/v_0^2}$ from $\vmin$ to $v_{\mathrm{esc}} - v_{\mathrm{E}}$, which with $y = v + v_{\mathrm{E}}$ is over $\vmin + v_{\mathrm{E}} < y < v_{\mathrm{esc}}$, and the constant $\ee^{-v_{\mathrm{esc}}^2/v_0^2}$ over an interval of length $2v_{\mathrm{E}}$. The Gaussian pieces combine by \cref{eq:D.1-erf}, the $\operatorname{erf}(v_{\mathrm{esc}}/v_0)$ terms cancel, and
\begin{equation}
  \eta(\vmin) = \frac{\pi v_0^2}{v_{\mathrm{E}}N_{\mathrm{esc}}}\left\{\frac{\sqrt{\pi}\,v_0}{2}\left[\operatorname{erf}\!\left(\frac{\vmin + v_{\mathrm{E}}}{v_0}\right) - \operatorname{erf}\!\left(\frac{\vmin - v_{\mathrm{E}}}{v_0}\right)\right] - 2v_{\mathrm{E}}\,\ee^{-v_{\mathrm{esc}}^2/v_0^2}\right\}.
  \label{eq:D.1-eta1}
\end{equation}

\emph{(ii) $v_{\mathrm{esc}} - v_{\mathrm{E}} \le \vmin < v_{\mathrm{esc}} + v_{\mathrm{E}}$.} Only the constant part of $e_{\mathrm{cut}}$ remains, over the length $v_{\mathrm{esc}} + v_{\mathrm{E}} - \vmin$, and
\begin{equation}
  \eta(\vmin) = \frac{\pi v_0^2}{v_{\mathrm{E}}N_{\mathrm{esc}}}\left\{\frac{\sqrt{\pi}\,v_0}{2}\left[\operatorname{erf}\!\left(\frac{v_{\mathrm{esc}}}{v_0}\right) - \operatorname{erf}\!\left(\frac{\vmin - v_{\mathrm{E}}}{v_0}\right)\right] - (v_{\mathrm{esc}} + v_{\mathrm{E}} - \vmin)\,\ee^{-v_{\mathrm{esc}}^2/v_0^2}\right\}.
  \label{eq:D.1-eta2}
\end{equation}

\emph{(iii) $\vmin \ge v_{\mathrm{esc}} + v_{\mathrm{E}}$.} No particle is fast enough: $\eta(\vmin) = 0$.

\emph{Checks.} At $\vmin = v_{\mathrm{esc}} - v_{\mathrm{E}}$ both \cref{eq:D.1-eta1,eq:D.1-eta2} give $(\sqrt{\pi}v_0/2)[\operatorname{erf}(v_{\mathrm{esc}}/v_0) - \operatorname{erf}((v_{\mathrm{esc}} - 2v_{\mathrm{E}})/v_0)] - 2v_{\mathrm{E}}\ee^{-v_{\mathrm{esc}}^2/v_0^2}$ inside the braces, so $\eta$ is continuous. At $\vmin = v_{\mathrm{esc}} + v_{\mathrm{E}}$ \cref{eq:D.1-eta2} vanishes, which matches branch (iii). At $\vmin = 0$ the two error functions in \cref{eq:D.1-eta1} are $\operatorname{erf}(v_{\mathrm{E}}/v_0)$ and $\operatorname{erf}(-v_{\mathrm{E}}/v_0) = -\operatorname{erf}(v_{\mathrm{E}}/v_0)$, and $\eta(0)$ is \cref{eq:1.2-inv}. Differentiating either branch with $\dd\operatorname{erf}(x)/\dd x = (2/\sqrt{\pi})\ee^{-x^2}$ gives $\dd\eta/\dd\vmin = -F_\chi(\vmin)/\vmin$, as the definition requires. Without the cut ($v_{\mathrm{esc}}\to\infty$, $N_{\mathrm{esc}}\to\pi^{3/2}v_0^3$) branch (i) becomes $[\operatorname{erf}((\vmin + v_{\mathrm{E}})/v_0) - \operatorname{erf}((\vmin - v_{\mathrm{E}})/v_0)]/(2v_{\mathrm{E}})$, the mean inverse speed of a shifted Maxwellian above $\vmin$.

